\subsection{The arithmetic family containing the prescribed key datum}
\label{subsec:ftb-arithmetic-key-family}

A key datum \(x=(A,L,\mathrm{rig},\theta,v)\) consists of an abelian
scheme of positive dimension \(g\), a rigidified symmetric principal
line bundle, a fibrewise nonzero section of that line bundle, and
an ordered Lie basis.  A key pullback includes all these data.
Its complex normal value is \(t(x)=\det(M)^4/c^8\), in the
normalization of Lemma~\ref{lem:ftb-key-normal-value-unique}.
Over a number field \(k\), a key datum is \emph{known} if one
\(t_0\in k\) gives its normal value under every embedding of \(k\)
in \(\mathbb C\), and for every corresponding key pullback.
A polarized variety is \emph{potentially known} if, after one
finite extension, every key datum on every further pullback of
that polarized variety is known.
Here the polarization includes the actual line bundle.

We separate two precise moduli inputs from the algebraic family
construction.  The first is a fine principal Siegel family \(F/K\)
at an odd level \(k\ge3\), with \(K\) a number field and
\(\zeta\in K\) a primitive \(k\)-th root.
It is smooth of relative dimension \(g(g+1)/2\), quasiprojective,
and has a universal principal polarization biextension and full
level.  It represents principal Siegel data over reduced bases
with dense characteristic-zero generic fibres, and separates
maps from local affine schemes giving the same Siegel datum.
The second input says that, for
\(\sigma\zeta=\exp(-2\pi i/k)\), any two complex points with
standard period-level data lie on the same connected component
of \(F(\mathbb C)\).
Their actual lower producers are recorded below; neither input
is replaced by an unnamed moduli or connectedness theorem.

\begin{srlemma}[A finite \'{e}tale torsion frame cover]
\label{lem:ftb-torsion-frame-cover}
For a Noetherian ring \(R\), an abelian scheme \(A/R\) of relative
dimension \(g\), and \(n\in R^\times\), there is a finite
\'{e}tale faithfully flat \(R\)-algebra \(R_1\) with
\[
             A[n]_{R_1}\simeq(\mathbb Z/n\mathbb Z)^{2g}_{R_1}.
\]
In particular the labelled sections list the full torsion,
without repetitions, after every algebraically closed field
pullback.
\end{srlemma}
\begin{proof}
The torsion scheme is finite \'{e}tale, and the geometric
\(d\)-torsion has \(d^{2g}\) elements for every \(d\mid n\).
These counts show that each geometric \(n\)-torsion group is
free of rank \(2g\).  Indeed, after choosing \(s<2g\) independent
generators, let \(m\) be the exponent of the quotient by their
span, and write \(n=mq\).  Multiplication by \(m\) has an image
of order \(q^{2g}\).  That image is generated by the images of
the chosen \(s\) generators: \(mx\) belongs to their span by
definition of \(m\), and their independence makes its
coordinates divisible by \(m\).
Thus \(q^{2g}\le q^s\), so \(q=1\).
Lift an element of order \(n\) in the quotient to obtain the
next independent generator.  Induction gives a frame.

In the product of \(2g\) copies of \(A[n]\), remove every
nonzero-coefficient linear-relation locus.  Each such locus
is open and closed, since the diagonal of a finite \'{e}tale
scheme is open and closed.  The remaining subscheme \(W\) is
finite \'{e}tale and is surjective by the geometric argument.
It is affine; let \(R_1\) be its coordinate ring.
Finite flatness and surjectivity give faithful flatness.
The tautological frame is an isomorphism on every geometric
fibre, hence an isomorphism of finite \'{e}tale group schemes.
\end{proof}

\begin{srlemma}[The actual symmetric principal square root]
\label{lem:ftb-symmetric-principal-root-cover}
Let \(R\) be a Noetherian \(\mathbb Q\)-algebra and \(P\) a
principal polarization biextension on an abelian scheme \(A/R\).
Put \(D=\Delta^*P\).
After a finite \'{e}tale faithfully flat change of base there
is a rigidified symmetric relatively ample line \(L\) and an
actual line-bundle isomorphism
\[
                             \lambda_L\simeq P.
 \tag{F1}\label{eq:ftb-family-actual-square-root}
\]
The full \(2\)-torsion is constant on this cover.
\end{srlemma}
\begin{proof}
Apply Lemma~\ref{lem:ftb-torsion-frame-cover} at \(n=2\), and
write \(C=(\mathbb Z/2\mathbb Z)^{2g}\) for the labelled
kernel, with basis \(h_i\).
Biextension additivity and symmetry give, on every test scheme,
\[
 D(x+h)=D(x)+D(h),\quad 2D(h)=0,\quad D(-x)=D(x)
                                                  \quad(2h=0).
 \tag{F2}\label{eq:ftb-diagonal-two-torsion-pic}
\]
Here \(D(x)=x^*[D]\) is a Picard class.
For example, if \(B(x,y)=(x,y)^*[P]\), then \(D(x)=B(x,x)\);
the cross terms are \(2B(x,h)=B(x,2h)=0\).

Choose generators \(\psi_i,\tau_i,j_0\) of, respectively,
\(D-t_{h_i}^*D+D(h_i)\), \(-2D(h_i)\), and \(D-[-1]^*D\).
Their existence is exactly \eqref{eq:ftb-diagonal-two-torsion-pic}.
For \(D^{\otimes2}\), set \(v_i=\psi_i^2\tau_i\) and \(j=j_0^2\).
Then
\begin{align}
 v_i\,t_{h_i}^*v_i&=a_i\in R^\times,\\
 v_i\,t_{h_i}^*v_j&=v_j\,t_{h_j}^*v_i,\\
 v_i\,t_{h_i}^*j&=j\,[-1]^*v_i.                  \tag{F3}
 \label{eq:ftb-diagonal-linearization-relations}
\end{align}
Here is the source's cancellation, rather than an assumption
that arbitrary theta-group lifts commute.
A quotient with zero formal Picard expression is a global unit
on \(A\), hence a base unit by \(\Gamma(A,\mathcal O_A)=R\).
Translations and inversion fix that quotient.
In particular \(\psi_i t_{h_i}^*\psi_i\tau_i\) is fixed by
\(t_{h_j}\), and
\[
 \frac{\psi_i t_{h_i}^*\psi_j}{\psi_j t_{h_j}^*\psi_i}
\]
is fixed by \(t_{h_i}\).
Multiply those two invariance identities and cancel to obtain
the second equality in \eqref{eq:ftb-diagonal-linearization-relations}.
For the third, use the invariance under inversion of
\(\psi_i t_{h_i}^*j_0/(j_0[-1]^*\psi_i)\), and the invariance
under \(t_{h_i}\) of \(j_0[-1]^*j_0\).
Their quotient gives the squared identity with its \(\tau_i\)
factors.  The first equality has zero formal Picard expression.

Adjoin \(b_i^2=a_i\) successively.
Each polynomial \(X^2-a_i\) is monic, its root is a unit, and
its derivative \(2b_i\) is a unit.
These extensions are finite free, \'{e}tale and faithfully flat.
The \(w_i=v_i/b_i\) are commuting involutions.
Their ordered products give generators \(u_c\), for \(c\in C\),
with the actual cocycle identity
\[
                           u_{c+c'}=u_c\,t_c^*u_{c'}.
 \tag{F4}\label{eq:ftb-family-linearization-cocycle}
\]
This follows by putting commuting involutions in basis order.
Their compatibility with \(j\) also persists.

The map \([2]:A\to A\) is finite \'{e}tale and surjective.
Two points with equal doubles differ \'{e}tale locally by a
unique labelled element of its kernel.
Thus \eqref{eq:ftb-family-linearization-cocycle} is a descent
cocycle on the actual fibre product.
The source constructs invariant local generators on each finite
orbit, shrinks the target until they are generators, and
descends their unit ratios.
They glue a line \(M\) with
\[
                                [2]^*M\simeq D^{\otimes2}.
 \tag{F5}\label{eq:ftb-family-double-pullback}
\]
The equality with \(j\) makes the isomorphism with \([-1]^*M\)
invariant under the kernel.  Its ratios descend by faithful
flatness, so \(M\) is symmetric.

To check its polarization, let \(N=\lambda_M-P\) on \(A\times A\).
It is trivial on both axes; the cube law makes it biadditive.
Evaluate \eqref{eq:ftb-family-double-pullback} on test points
to obtain \(N(2x,2y)=0\), and hence \(N(4x,y)=0\).
One must not divide this equality in the Picard group.
Instead use the closed triviality locus for this line family
in the first factor.  It contains the image of \([4]\).
A closed immersion through which a faithfully flat surjection
factors is an isomorphism.
Consequently the line is locally trivial in that factor.
Its trivialization on the other axis and the equality of
global functions with those on the base give a global
trivialization.  This proves \(\lambda_M\simeq P\) as a line
isomorphism.

Replace \(M\) by \(M-e^*M\).
It is rigidified and symmetric, with the same polarization.
The cube law, now using that actual polarization identity,
identifies its square with \(D\).
It is relatively ample because \(D\) is; it is principal
because its polarization is \(P\).
All torsion labels survive these changes of base.
\end{proof}

\begin{srlemma}[Keeping both points on one integral component]
\label{lem:ftb-two-point-etale-integral-component}
Let \(K,K_1\) be number fields with \(K_1\hookrightarrow\mathbb C\).
Let \(R\) be a smooth finitely generated \(K\)-algebra and
\(R\to R_1\) finite \'{e}tale faithfully flat.
Suppose \(p_1:R_1\to\mathbb C\) and \(q:R\to K_1\) are
\(K\)-maps, and \(p=p_1|_R\).
Assume \(p\) and \(q\) agree on all idempotents after extending
scalars to \(\mathbb C\).
There exist a number field \(k\supset K_1\), a smooth
geometrically integral finitely generated \(k\)-algebra \(S\),
a map \(R_1\to S\), a complex point \(s\), and a \(k\)-point
\(s_0\), with
\[
              s|_{R_1}=p_1,\qquad
              s_0|_R=(K_1\hookrightarrow k)q .
 \tag{F6}\label{eq:ftb-two-point-component-maps}
\]
\end{srlemma}
\begin{proof}
Inside \(\mathbb C\), let \(\overline K\) be an algebraic closure
of \(K\).  The ring \(\overline K\otimes_KR_1\) is regular.
Its finitely many irreducible components are disjoint, and
therefore integral open and closed components.
Choose the one containing \(p_1\), with idempotent \(e\).
Its finitely many coefficients descend to a finite extension
\(k_0/K\).
The corresponding localization is geometrically integral:
its extension to \(\overline K\) is a domain, and the
algebraically closed field test gives geometric integrality.

The image of this component in
\(\operatorname{Spec}(\overline K\otimes_KR)\) is open and
closed, by flatness with finite presentation and by finiteness.
It contains \(p\).  Agreement on idempotents puts \(q\) in
that image, so \(q\) has a lift on this same component.
The finitely many coordinates of the lift are algebraic over
\(K_1\), since \(R_1/R\) is finite.
Adjoin them and \(k_0\) to \(K_1\), obtaining \(k\).
The base change of the selected component is \(\operatorname{Spec}S\).
The selected lift gives \(s_0\), while \(p_1\) gives \(s\).
Both identities in \eqref{eq:ftb-two-point-component-maps}
are built into the construction.
\end{proof}

\begin{srlemma}[Knownness and the actual polarization ambiguity]
\label{lem:ftb-known-same-polarization}
A polarized abelian variety over a number field having the
same principal polarization biextension as a known one is
potentially known.
\end{srlemma}
\begin{proof}
Over the algebraic closure, identify the two abelian varieties.
The quotient of their symmetric principal lines is in
\(\operatorname{Pic}^0[2]\).
The principal polarization identifies it with a \(2\)-torsion
point, so the lines differ by translation by that point.
The point and a line isomorphism descend to one finite
extension, by adjoining their finitely many coordinates.

The line normalization is essential.
For the translated key datum its cube/rigidification data
supplies a unit \(\nu\).
In an old normal form with \(c=1,M=1\), write the point as
\((m+\tau n)/2\).
The half-period theta formula gives the new characteristic
and scalar \(c'\) with
\[
                               (c')^8\nu^2=1.
\]
Explicitly, let \(\lambda\) be the value of the translation
generator at that half period, and let \(d\) be the scalar
in the theta translation formula.
The normalized doubling generator has value \(1\).
Evaluation at the half period and its double gives
\(\lambda^8=\nu^2d^8\), and \(c'=d/\lambda\).
Thus translation multiplies the normal value by \(\nu^2\).
All these identities are over the same finite extension,
and are preserved by every field embedding.

For arbitrary frames the exact change rule is
\[
                    t(a\theta,Bv)=\det(B)^4a^{-8}t(\theta,v).
 \tag{F7}\label{eq:ftb-known-frame-change}
\]
Its multiplier belongs to the field of definition.
Hence every key frame on every further pullback is known.
\end{proof}

\begin{srproposition}[The arithmetic connected-family join]
\label{prop:ftb-connfam-arithmetic-join}
Under the two stated moduli inputs, with their actual period,
level and polarization pullback comparisons, every complex key
datum satisfies \eqref{eq:ftb-connfam}.
\end{srproposition}
\begin{proof}
Fix \(x/\mathbb C\) of dimension \(g\).
The first input gives \(F/K,k,\zeta,\sigma\).
Represent the period matrix and standard level of \(x\)
by \(p:\operatorname{Spec}\mathbb C\to F\).
For the known point use the \(g\)-fold product of
\(E:y^2=x^3-x\), of discriminant \(64\).
The elliptic normalization and product identity give normal
value \(64^g\).
After a finite number-field extension \(K_1/K\), all its
\(k\)-torsion and the compatible level are defined.
Represent it by \(q:\operatorname{Spec}K_1\to F\).
The second input puts its complex pullback and \(p\) on the
same component.

Choose a single affine open \(\operatorname{Spec}R\) containing
both points.  Quasiprojectivity gives it by a hypersurface
avoiding their finite set of images.
It is smooth over \(K\).
The required agreement on idempotents of
\(\mathbb C\otimes_KR\) follows from the ambient component
condition: a smooth Noetherian complex scheme is regular,
so its irreducible components are disjoint; the intersection
of any one with the affine open remains irreducible.
Thus the two points lie on the same component of this
affine open as well.

Apply Lemma~\ref{lem:ftb-symmetric-principal-root-cover} to the
universal biextension on \(R\), obtaining \(R_1,Q\), and lift
\(p\) to \(p_1\).
The fibre of \(Q\) and the line of \(x\) have the same
biextension, so their discrepancy is a \(2\)-torsion translate.
The full \(2\)-torsion is constant over \(R_1\), and the
required point therefore lifts to a section.
Translate \(Q\) by that section and rigidify by its value
there.  The resulting \(Q'\) has the actual line of \(x\)
at \(p_1\), and its principal biextension is unchanged.

Lemma~\ref{lem:ftb-two-point-etale-integral-component} gives
\(k,S,s,s_0\).
The pullback of \(Q'\) has the prescribed polarized fibre
at \(s\); its fibre at \(s_0\) has the biextension of the
known product.
Lemma~\ref{lem:ftb-known-same-polarization} makes that fibre
potentially known.
Make its finite constant-field extension now.
Geometric integrality preserves a smooth integral base;
the complex embedding extends, and both points lift.

The theta sections form an invertible module, by principal
polarization cohomology and base change; the Lie module is
projective of rank \(g\).
At the two primes defined by \(s,s_0\), both are free over
the semilocal ring.
Choose their frames there, clear the finitely many
denominators, and localize by one element outside both
primes.  This gives a key datum \(y\) while retaining both
points.
Before comparing frames, transport the rigidification of \(x\)
through the actual polarized line-bundle pullback isomorphism.
The source's key-on-a-fixed-polarization comparison identifies
that transported rigidification with the chosen one; its line
isomorphism is included in this comparison.  The remaining
parameters are a theta scalar and a Lie matrix.
To impose exactly the specified frames of \(x\), use
\[
                  S'=S[X,X_{ij},(X\det(X_{ij}))^{-1}].
 \tag{F8}\label{eq:ftb-universal-key-frame-torsor}
\]
Multiply the theta section by \(X\), and the Lie frame by
the invertible matrix \((X_{ij})\).
At the complex point take the scalar and matrix carrying
the existing frames to those of \(x\); at the known point
take \(1,I_g\).
The new base is smooth and geometrically integral.
Its complex fibre is the exact key datum \(x\), while the
known fibre is unchanged.
Knownness gives one \(t_0\in k\) at every embedding and for
every corresponding key pullback.
These are precisely the clauses of \eqref{eq:ftb-connfam}.
\end{proof}

\paragraph{Frozen source concordance and the separate lower producers.}
The prefix is
\srcid{src/Development/AbelianVarietyTheory/Proofs/}.
The torsion frame is
\srcid{ThetaLawProofs/ExistsFiniteEtaleCoverConstantTorsion.lean}
(\srcid{f039af8cf45f309722d9c38fdb77e32c50665d19}).
The actual symmetric root uses, in \srcid{ThetaLawProofs},
\srcid{ExistsSymmetricLinearisationSquareDiagonalBundle.lean}
(\srcid{994c2070f67e20bdb9ce54c8efabc9c82511efb9}),
\srcid{ExistsDescentSymmetricLinearisedBundle.lean}
(\srcid{b4d02b09f887aa1bc7e23314b3304db2ee82f9d8}),
\srcid{LamIsoOfDoublingPullbackEqDiagonalSq.lean}
(\srcid{87e1d8620f1761e05f668291e221f8b66f0c781d}),
\srcid{ExistsPrincipalPolarizationSymmetricSquareRoot.lean}
(\srcid{481ca29a56b6d87791856ae5d24b605e3a0cc2d0}),
and \srcid{ExistsSymmetricSqrtPrincipalSiegelFiniteEtale.lean}
(\srcid{ad982d88a3039d9157db361c7eb47c1d2158141f}).
The two-point component argument is
\srcid{GeomIntegralOfEtaleFaithfullyFlatCover.lean}
(\srcid{2c3cdb5d68e977f5fcbac8949e3f558d7ceeb1ce}),
with \srcid{ExistsFiniteExtGeomIntegralLocalization.lean}
(\srcid{ce4a22acec9b2c4be3162269c3732e106c48ecd4}).
Knownness is
\srcid{PotentiallyKnownOfSameLambdaBundle.lean}
(\srcid{46b881f6de2bcd933fdf0af3f175bf70499ce687}),
\srcid{ThetaFrameValueKnownOfTwoTorsionTranslate.lean}
(\srcid{5cbd8eaa06fc6135efd5dcd411fd1f942367eb26}),
\srcid{ThetaNormalFormTranslateHalfPeriod.lean}
(\srcid{72dd6206484af0e01b4c33ad7eecbbb917c36da1}),
\srcid{KeyNormalFormDetRatioTwoTorsionTranslate.lean}
(\srcid{3429a6c1eec5af7b9062c8490ea52e893ec1aa9d}),
and \srcid{ThetaFrameValueKnownOfPolarizedPullback.lean}
(\srcid{391e84ccf8fbd2fe2dfae08549198056345d0a0e}).
The final frame construction is
\srcid{ExistsKeyDataNearTwoPrimes.lean}
(\srcid{b08147dbe77819b1d425e3345501c1fb16b9a1eb})
and \srcid{ExistsThetaKeyFamilyUniversalFrames.lean}
(\srcid{8950e79c3f27fed92a98df44b05825d3e1cfdb39}).
The joins are
\srcid{PolarizedFamilyOfThetaUniformization.lean}
(\srcid{6e9fb2cd94e212f61adaaa66272037b6e4195b01})
and \srcid{KeyDatumConnectedFamilyOfPolarizedArithmeticFami.lean}
(\srcid{3414212fc4b9f57646d6e2b692d061f799a58244}).

The first still separate moduli producer is
\srcid{ModuliProofs/NonemptySiegelFamilyOverArithmeticBase.lean}
(\srcid{c8933879c627986eb6412db2774dfd3f8d998c8b}).
It constructs the fine family through the actual eighth-power
sharp theta closure; generic dimension, special tangent
dimension, universal mapping and principal-Siegel recovery
are substantive lower lemmas.
Its number-field specialization is
\srcid{ThetaLawProofs/ExistsSiegelFamilyOverNumberField.lean}
(\srcid{a87124f2cc26de113f66ea97d5641108f3ff26c9}).
The second input is
\srcid{SiegelFamilyStandardLevelPointsSameComponent.lean}
(\srcid{26a36ef9ff5a0ee06e4f6ba603ea75f78fbaa829}).
Its local producer is
\srcid{SiegelFamilyNearbyPeriodsSameComponent.lean}
(\srcid{0908e63fa80faff82352cd114d3ce644623a521a}),
using actual relative periods, level/Weil compatibility,
uniqueness at a fixed period matrix, and openness of the
locally injective period chart.
These and the fixed-fibre complex realization remain
explicit until their own translations are joined.
